% Curso intensivo del Capítulo 11 de la tesis (completitud de la deducción
% natural proposicional) y del Apéndice A (corrección, consistencia, adecuación).
% Escrito en matemática de papel, sin Lean. Fuente:
%   tesis-mate/documento/capitulos/4_practica.tex
%   tesis-mate/documento/capitulos/apendice_correccion.tex
% Compilar dos veces:  lualatex CursoCompletitud.tex
\documentclass[11pt,a4paper]{article}
\usepackage[a4paper,margin=2.3cm,headheight=14pt]{geometry}
\usepackage{amsmath,mathtools}
\usepackage{fontspec}
\usepackage{unicode-math}
\usepackage[spanish,es-noshorthands,es-tabla]{babel}
\setmainfont{TeX Gyre Pagella}
\setmathfont{TeX Gyre Pagella Math}
\setsansfont{DejaVu Sans}[Scale=MatchLowercase]
\setmonofont{DejaVu Sans Mono}[Scale=MatchLowercase]
\usepackage{xcolor}
\usepackage[most]{tcolorbox}
\usepackage{booktabs,tabularx,array,longtable}
\usepackage{bussproofs}
\usepackage{tikz}
\usetikzlibrary{arrows.meta,positioning,trees}
\usepackage{enumitem}
\usepackage{fancyhdr}
\usepackage[hidelinks]{hyperref}
\usepackage{microtype}

\definecolor{fondores}{HTML}{F4F7FB}
\definecolor{borderes}{HTML}{3F6FA8}
\definecolor{fondoesquema}{HTML}{EEF5EE}
\definecolor{bordeesquema}{HTML}{5E9A62}
\definecolor{fondonota}{HTML}{F1EEF8}
\definecolor{bordenota}{HTML}{8C7DB8}
\definecolor{fondoejemplo}{HTML}{F7F8FA}
\definecolor{bordeejemplo}{HTML}{B8C2CF}
\definecolor{fondotrampa}{HTML}{FFF4EC}
\definecolor{bordetrampa}{HTML}{D98A4E}
\definecolor{fondoclave}{HTML}{FFF9E6}
\definecolor{bordeclave}{HTML}{C9A227}

\newtcolorbox{resultado}[1]{enhanced, breakable, colback=fondores,
  colframe=borderes, boxrule=0.6pt, arc=3pt, colbacktitle=borderes,
  coltitle=white, fonttitle=\sffamily\bfseries, title={#1},
  before skip=10pt, after skip=10pt}
\newtcolorbox{esquema}[1]{enhanced, breakable, colback=fondoesquema,
  colframe=bordeesquema, boxrule=0.6pt, arc=3pt,
  colbacktitle=bordeesquema, coltitle=white, fonttitle=\sffamily\bfseries,
  title={#1}, before skip=10pt, after skip=10pt}
\newtcolorbox{ejemplo}[1]{enhanced, breakable, colback=fondoejemplo,
  colframe=bordeejemplo, boxrule=0.5pt, arc=2pt, colbacktitle=bordeejemplo!40,
  coltitle=black, fonttitle=\sffamily\small\bfseries, title={#1},
  before skip=6pt, after skip=8pt}
\newtcolorbox{nota}[1][Nota]{enhanced, breakable, colback=fondonota,
  colframe=bordenota, boxrule=0.5pt, arc=3pt, colbacktitle=bordenota,
  coltitle=white, fonttitle=\sffamily\bfseries\small, title={#1},
  before skip=8pt, after skip=8pt}
\newtcolorbox{trampa}[1][Lo que pueden preguntar]{enhanced, breakable,
  colback=fondotrampa, colframe=bordetrampa, boxrule=0.5pt, arc=3pt,
  colbacktitle=bordetrampa, coltitle=white, fonttitle=\sffamily\bfseries\small,
  title={#1}, before skip=8pt, after skip=8pt}
\newtcolorbox{clave}{enhanced, colback=fondoclave, colframe=bordeclave,
  boxrule=0.8pt, arc=3pt, left=8pt, right=8pt, top=6pt, bottom=6pt,
  before skip=10pt, after skip=10pt}

\newcommand{\campo}[1]{\par\smallskip\noindent\textbf{#1}\enspace}
\newcommand{\enunciado}{\campo{Enunciado.}}
\newcommand{\palabras}{\campo{En palabras.}}
\newcommand{\sirve}{\campo{Para qué sirve.}}
\newcommand{\idea}{\campo{Idea.}}
\newcommand{\prueba}{\campo{Demostración.}}

% Notación de la tesis
\newcommand{\Provnd}{\vdash_{\mathrm{ND}}}
\newcommand{\Imp}{\rightarrow}
\newcommand{\Not}{\neg}
\newcommand{\sem}[1]{[\![#1]\!]}
\newcommand{\Taut}{\operatorname{Taut}}
\newcommand{\atomos}{\operatorname{atomos}}
\newcommand{\lit}{\operatorname{literal}}
\newcommand{\Nodup}{\operatorname{Nodup}}
\newcommand{\dedup}{\operatorname{dedup}}
\newcommand{\V}{\mathsf{V}}
\newcommand{\F}{\mathsf{F}}
\newcommand{\lean}[1]{{\footnotesize\texttt{#1}}}

\pagestyle{fancy}
\fancyhf{}
\fancyhead[L]{\small\sffamily Capítulo 11: el teorema de completitud}
\fancyhead[R]{\small\sffamily\nouppercase{\leftmark}}
\fancyfoot[C]{\small\thepage}
\renewcommand{\headrulewidth}{0.4pt}

\setlist{itemsep=2pt, topsep=4pt}
\setlength{\parskip}{4pt}
\setlength{\parindent}{0pt}
\EnableBpAbbreviations

\title{\textbf{Curso intensivo del Capítulo 11}\\[4pt]
  \Large El teorema de completitud de la deducción natural\\[2pt]
  \large Qué se demuestra, cómo se demuestra y cómo explicarlo}
\author{Material de apoyo para la defensa}
\date{}

\begin{document}
\maketitle
\thispagestyle{empty}

\begin{abstract}
\noindent
Este documento explica el Capítulo 11 de la tesis y su complemento, el
Apéndice A, en matemática de papel. No usa código. Cada resultado aparece en una
ficha con la misma estructura: enunciado, qué dice en palabras, para qué sirve,
la idea en una frase y la demostración paso a paso. Al final hay un ejemplo
resuelto de principio a fin, un análisis de dónde interviene la lógica clásica y
las preguntas más probables del tribunal con sus respuestas. La numeración
(Lema 11.15, Teorema A.3, etc.) es la de la versión compilada de la tesis.
\end{abstract}

\tableofcontents
\newpage

%=====================================================================
\section{El capítulo en una página}
%=====================================================================

\begin{clave}
\centering
\textbf{Teorema principal (Teorema 11.18).}\quad
Toda tautología es derivable sin hipótesis:
\[
\Taut(\varphi)\quad\Longrightarrow\quad\varnothing\Provnd\varphi .
\]
\end{clave}

\textbf{Qué significa.} «\(\varphi\) es verdadera en todas las filas de su
tabla de verdad» (un hecho \emph{semántico}) implica «existe una demostración
formal de \(\varphi\) en deducción natural» (un hecho \emph{sintáctico}). Es
decir, las reglas de la deducción natural bastan para demostrar todo lo que es
lógicamente verdadero. La dirección contraria, \(\varnothing\Provnd\varphi
\Rightarrow\Taut(\varphi)\), es la \emph{corrección} y está en el Apéndice A.
Las dos juntas forman la \emph{adecuación} (Teorema A.6).

\textbf{Cómo se demuestra.} Con el argumento de Kalmár, en dos etapas:

\begin{enumerate}
  \item \textbf{Una derivación por fila (Lema 11.15).} Se fija una fila de la
    tabla de verdad, es decir, una valuación \(v\). Se toman como hipótesis
    los \emph{literales} de esa fila: \(p\) si \(v(p)=\V\) y \(\Not p\) si
    \(v(p)=\F\). Con esas hipótesis se deriva \(\varphi\) si es verdadera en la
    fila, o \(\Not\varphi\) si es falsa.
  \item \textbf{Juntar las filas (Lema 11.17).} Si \(\varphi\) es tautología,
    hay una derivación de \(\varphi\) en \emph{cada} fila. Se eliminan las
    hipótesis átomo por átomo: de «\(\varphi\) se deriva suponiendo \(p\)» y
    «\(\varphi\) se deriva suponiendo \(\Not p\)» se obtiene «\(\varphi\) se
    deriva sin suponer nada sobre \(p\)» (Lema 11.11). Al acabar con todos los
    átomos no queda ninguna hipótesis.
\end{enumerate}

\subsection*{Mapa de dependencias}

Cada flecha va de un resultado al que lo usa. Las cajas punteadas son las del
Apéndice A (rama de la corrección).

\begin{center}
\begin{tikzpicture}[font=\footnotesize,
  res/.style={draw, rounded corners, fill=blue!6, align=center, inner sep=3pt, text width=26mm},
  ap/.style={draw, dashed, rounded corners, fill=gray!5, align=center, inner sep=3pt, text width=26mm},
  dep/.style={-{Stealth}, thick, gray!60!black}]
\node[res] (weak)    at (0,0)      {Lema 11.9\\debilitamiento};
\node[res] (litmem)  at (3.3,0)    {Lema 11.14\\pertenencia del literal};
\node[res] (congr)   at (6.6,0)    {Lema 11.16\\congruencia};
\node[res] (dni)     at (0,-1.7)   {Lema 11.10\\doble negación};
\node[res] (bycases) at (6.6,-1.7) {Lema 11.11\\análisis por casos};
\node[res] (kalmar)  at (1.65,-3.4){Lema 11.15\\Kalmár};
\node[res] (disch)   at (6.6,-3.4) {Lema 11.17\\descarga};
\node[res, fill=yellow!20] (compl) at (4.1,-5.1) {\textbf{Teorema 11.18}\\completitud};
\node[ap] (sound)    at (10.2,-1.7) {Teorema A.3\\corrección};
\node[ap] (taut)     at (10.2,-3.4) {Corolario A.4\\teoremas = tautologías};
\node[ap] (cons)     at (12.0,-5.1) {Corolario A.5\\consistencia};
\node[ap] (adeq)     at (8.3,-5.1)  {Teorema A.6\\adecuación};
\draw[dep] (weak) -- (dni);
\draw[dep] (weak) -- (bycases);
\draw[dep] (weak) -- (kalmar);
\draw[dep] (litmem) -- (kalmar);
\draw[dep] (dni) -- (kalmar);
\draw[dep] (bycases) -- (disch);
\draw[dep] (congr.east) to[bend left=35] (disch.east);
\draw[dep] (kalmar) -- (compl);
\draw[dep] (disch) -- (compl);
\draw[dep] (sound) -- (taut);
\draw[dep] (taut) -- (cons);
\draw[dep] (taut) -- (adeq);
\draw[dep] (compl) -- (adeq);
\end{tikzpicture}
\end{center}


\subsection*{Todos los resultados de un vistazo}

{\small
\begin{longtable}{@{}p{0.14\textwidth}p{0.47\textwidth}p{0.33\textwidth}@{}}
\toprule
\textbf{Resultado} & \textbf{Enunciado} & \textbf{Técnica}\\
\midrule
\endhead
Lema 11.9 & \(\Gamma\Provnd\varphi\) y \(\Gamma\subseteq\Delta\) \(\Rightarrow\) \(\Delta\Provnd\varphi\) & inducción sobre la derivación, con \(\Delta\) generalizado\\
Lema 11.10 & \(\Gamma\Provnd\varphi\Rightarrow\Gamma\Provnd\Not\Not\varphi\) & construir una derivación\\
Lema 11.11 & \(\Gamma,\varphi\Provnd\chi\) y \(\Gamma,\Not\varphi\Provnd\chi\) \(\Rightarrow\) \(\Gamma\Provnd\chi\) & construir una derivación con RAA\\
Lema 11.14 & \(p\in A\Rightarrow\lit_v(p)\in\Gamma_v(A)\) & inducción sobre la lista\\
Lema 11.15 & Kalmár: cada fila deriva \(\varphi\) o \(\Not\varphi\) & inducción sobre \(\varphi\), enunciado bilateral\\
Lema 11.16 & valuaciones iguales en \(A\) dan \(\Gamma_v(A)=\Gamma_w(A)\) & casos e inducción sobre la lista\\
Lema 11.17 & \(\forall v,\ \Gamma_v(A)\Provnd\varphi\) \(\Rightarrow\) \(\varnothing\Provnd\varphi\) & inducción sobre la lista, modificando \(v\)\\
Teorema 11.18 & \(\Taut(\varphi)\Rightarrow\varnothing\Provnd\varphi\) & composición: Kalmár + descarga\\
\midrule
Lema A.1 & \(\sem{\bot}_v\) es falsa & cálculo directo\\
Lema A.2 & \(\sem\chi_v\) y \(v\models\Gamma\) \(\Rightarrow\) \(v\models\Gamma\cup\{\chi\}\) & casos sobre la pertenencia\\
Teorema A.3 & corrección: \(\Gamma\Provnd\varphi\Rightarrow(v\models\Gamma\Rightarrow\sem\varphi_v)\) & inducción sobre la derivación, con \(v\) generalizada\\
Corolario A.4 & \(\varnothing\Provnd\varphi\Rightarrow\Taut(\varphi)\) & caso \(\Gamma=\varnothing\)\\
Corolario A.5 & \(\varnothing\not\Provnd\bot\) & contradicción con A.1\\
Teorema A.6 & \(\Taut(\varphi)\iff\varnothing\Provnd\varphi\) & 11.18 y A.4\\
\bottomrule
\end{longtable}
}

%=====================================================================
\section{El escenario: definiciones}
%=====================================================================

\subsection{Fórmulas y verdad}

\begin{itemize}
  \item \textbf{Fórmulas (Def.~11.3).} Solo hay tres formas:
    \[\varphi::=p\ \mid\ \Not\varphi\ \mid\ (\varphi\Imp\psi),\]
    con \(p\) un átomo (una variable proposicional). Los demás conectivos se
    escriben como abreviaturas (Sección~\ref{sec:base}).
  \item \textbf{Valuación y evaluación (Def.~11.4).} Una valuación \(v\)
    asigna \(\V\) o \(\F\) a cada átomo. El valor de una fórmula se calcula por
    recursión:
    \[
    \sem{p}_v=v(p),\qquad
    \sem{\Not\varphi}_v=\text{no }\sem{\varphi}_v,\qquad
    \sem{\varphi\Imp\psi}_v=\big(\sem{\varphi}_v\Rightarrow\sem{\psi}_v\big).
    \]
  \item \textbf{Tautología (Def.~11.5).} \(\Taut(\varphi)\) significa que
    \(\sem\varphi_v\) es verdadera para \emph{toda} \(v\).
  \item \textbf{Satisfacción de un contexto (Def.~11.6).} \(v\models\Gamma\)
    significa que \(v\) hace verdaderas todas las fórmulas de \(\Gamma\).
  \item \textbf{La fórmula falsa (Def.~11.8).} No hay \(\bot\) primitivo. Se
    fija un átomo \(p_0\) y se define
    \[\bot:=\Not(p_0\Imp p_0),\]
    que es falsa en toda valuación, cualquiera que sea \(v(p_0)\) (Lema A.1).
\end{itemize}

\subsection{\texorpdfstring{El cálculo \(\Provnd\): siete reglas}{El cálculo ND: siete reglas}}

El juicio \(\Gamma\Provnd\varphi\) («\(\varphi\) se deriva de las hipótesis
\(\Gamma\)») es lo que generan estas siete reglas, y nada más:

\begin{center}
\AxiomC{\(\varphi\in\Gamma\)}\RightLabel{\scriptsize hip.}\UnaryInfC{\(\Gamma\Provnd\varphi\)}\DisplayProof
\qquad
\AxiomC{\(\Gamma,\varphi\Provnd\psi\)}\RightLabel{\scriptsize\(\Imp\)I}\UnaryInfC{\(\Gamma\Provnd\varphi\Imp\psi\)}\DisplayProof
\qquad
\AxiomC{\(\Gamma\Provnd\varphi\Imp\psi\)}\AxiomC{\(\Gamma\Provnd\varphi\)}\RightLabel{\scriptsize\(\Imp\)E}\BinaryInfC{\(\Gamma\Provnd\psi\)}\DisplayProof
\\[10pt]
\AxiomC{\(\Gamma,\varphi\Provnd\bot\)}\RightLabel{\scriptsize\(\Not\)I}\UnaryInfC{\(\Gamma\Provnd\Not\varphi\)}\DisplayProof
\qquad
\AxiomC{\(\Gamma\Provnd\Not\varphi\)}\AxiomC{\(\Gamma\Provnd\varphi\)}\RightLabel{\scriptsize\(\Not\)E}\BinaryInfC{\(\Gamma\Provnd\bot\)}\DisplayProof
\qquad
\AxiomC{\(\Gamma\Provnd\bot\)}\RightLabel{\scriptsize\(\bot\)E}\UnaryInfC{\(\Gamma\Provnd\varphi\)}\DisplayProof
\qquad
\AxiomC{\(\Gamma,\Not\varphi\Provnd\bot\)}\RightLabel{\scriptsize RAA}\UnaryInfC{\(\Gamma\Provnd\varphi\)}\DisplayProof
\end{center}

Las seis primeras forman el núcleo intuicionista. \textbf{RAA} (reducción al
absurdo clásica) es la única regla clásica, y es exactamente la que hace
falta para la completitud respecto de las tablas de verdad.

\subsection{Dos convenios que conviene saber explicar}

\begin{enumerate}
  \item \textbf{Los contextos son conjuntos} (§11.3.2). \(\Gamma,\varphi\)
    significa \(\Gamma\cup\{\varphi\}\). Por eso el orden y las repeticiones
    de hipótesis no importan (intercambio y contracción vienen gratis), pero
    \emph{añadir} hipótesis no viene gratis: el debilitamiento es un teorema
    sobre derivaciones (Lema 11.9).
  \item \textbf{La negación es primitiva y \(\bot\) es la abreviatura}
    (§11.3.2). Es al revés de lo habitual (\(\Not\varphi:=\varphi\to\bot\)).
    Por eso \(\Not\)I y \(\Not\)E son reglas propias y no casos de \(\Imp\)I e
    \(\Imp\)E.
\end{enumerate}

\subsection{Los objetos del argumento de Kalmár}

\begin{itemize}
  \item \textbf{Átomos de una fórmula (Def.~11.12):} la lista, con
    repeticiones, de los átomos que aparecen.
    \(\atomos(p)=[p]\), \(\atomos(\Not\varphi)=\atomos(\varphi)\),
    \(\atomos(\varphi\Imp\psi)=\atomos(\varphi)\mathbin{+\!\!+}\atomos(\psi)\).
  \item \textbf{Literal (Def.~11.13):} \(\lit_v(p)=p\) si \(v(p)=\V\) y
    \(\lit_v(p)=\Not p\) si \(v(p)=\F\). Es «lo que la fila dice de \(p\)».
  \item \textbf{Contexto de literales (Def.~11.13):}
    \(\Gamma_v(A)=\{\lit_v(p):p\in A\}\), con
    \(\Gamma_v([\,])=\varnothing\) y
    \(\Gamma_v(p::r)=\{\lit_v(p)\}\cup\Gamma_v(r)\).
  \item \textbf{Cobertura:} \(A\) \emph{cubre} a \(\varphi\) si
    \(\atomos(\varphi)\subseteq A\) (cada átomo de \(\varphi\) está en \(A\)).
  \item \textbf{Sin repeticiones y \(\dedup\) (§11.5.4--11.5.5):}
    \(\Nodup(A)\) dice que ningún átomo aparece dos veces. \(\dedup\) borra
    las apariciones repetidas, conservando la última:
    \[\dedup([p,q,p])=[q,p].\]
\end{itemize}

\begin{ejemplo}{Un contexto de literales}
Con \(A=[q,p]\) y la fila \(v(p)=\V\), \(v(q)=\F\):
\[
\Gamma_v([q,p])=\{\lit_v(q)\}\cup\{\lit_v(p)\}=\{\Not q,\ p\}.
\]
Hay una fila por valuación de los átomos de \(A\), es decir,
\(2^{2}=4\) contextos: \(\{q,p\}\), \(\{\Not q,p\}\), \(\{q,\Not p\}\) y
\(\{\Not q,\Not p\}\).
\end{ejemplo}

\subsection{\texorpdfstring{Por qué basta con \(\Not\) y \(\Imp\)}{Por qué basta con negación e implicación}}\label{sec:base}

La tesis justifica la base \(\{\Not,\Imp\}\) en §11.2.2 con una traducción
\(T\) del lenguaje completo (\(\wedge,\vee,\leftrightarrow,\top,\bot\)) al
fragmento:
\[
T(\varphi\wedge\psi)=\Not(T\varphi\Imp\Not T\psi),\qquad
T(\varphi\vee\psi)=\Not T\varphi\Imp T\psi,\qquad
T(\top)=p_0\Imp p_0,\qquad T(\bot)=\Not(p_0\Imp p_0),
\]
y \(T\) deja igual átomos, \(\Not\) e \(\Imp\). Por inducción estructural,
\(\sem{T(\theta)}_v=\sem{\theta}_v\) para toda \(v\), así que
\[
\models\theta\iff\models T(\theta)\iff\varnothing\Provnd T(\theta).
\]
La segunda equivalencia es la adecuación (Teorema A.6).

\begin{nota}[Dos precisiones que la tesis hace]
\begin{itemize}
  \item La preservación de \(T\) se justifica matemáticamente en la tesis,
    pero no está verificada en la formalización. Lo verificado es la
    adecuación para fórmulas de \(\{\Not,\Imp\}\).
  \item La traducción es correcta para la lógica \emph{clásica}, no para la
    intuicionista: \(T(p\vee\Not p)=\Not p\Imp\Not p\) es demostrable incluso
    sin RAA, mientras que \(p\vee\Not p\) no lo es intuicionistamente.
\end{itemize}
\end{nota}

%=====================================================================
\section{Cómo leer las demostraciones de la tesis}
%=====================================================================

\subsection{Dos niveles que no hay que mezclar}

La Nota 11.1 de la tesis distingue:
\begin{center}
\renewcommand{\arraystretch}{1.3}
\begin{tabularx}{\linewidth}{@{}p{0.2\linewidth}XX@{}}
\toprule
& \textbf{Nivel objeto} & \textbf{Nivel meta}\\
\midrule
Qué es & el cálculo \(\Provnd\) que se estudia & el razonamiento \emph{sobre} el cálculo\\
Negación & \(\Not\) (un símbolo de fórmula) & «no» (una afirmación)\\
Falsedad & \(\bot:=\Not(p_0\Imp p_0)\), una fórmula & «Falso», una proposición\\
Absurdo clásico & regla RAA del cálculo & razonar por contradicción\\
Ejemplo & \(\Gamma,\Not\varphi\Provnd\bot\Rightarrow\Gamma\Provnd\varphi\) & «si \(\sem\varphi_v\) fuera falsa, \dots; luego es verdadera»\\
\bottomrule
\end{tabularx}
\end{center}

Una frase como «por RAA» puede referirse a cualquiera de los dos niveles. Al
explicar una prueba conviene decir siempre en cuál se está.

\subsection{Las reglas R1--R17}

La tesis escribe cada paso del razonamiento meta citando una regla de la
Nota~11.2. Este curso usa las frases habituales; la equivalencia es:
\begin{center}
{\small
\begin{tabular}{@{}lll@{}}
\toprule
\textbf{Regla} & \textbf{Nombre} & \textbf{Frase habitual}\\
\midrule
R1, R2 & \(\Rightarrow\)-intro / elim & «supongamos \(P\) \dots\ luego \(P\Rightarrow Q\)»; «como \(P\Rightarrow Q\) y \(P\), \(Q\)»\\
R3, R4 & \(\wedge\)-intro / elim & «se tienen ambas»; «en particular»\\
R5, R6 & \(\vee\)-intro / casos & «o bien \dots\ o bien \dots; en ambos casos»\\
R7, R8 & \(\iff\)-intro / elim & «se prueban las dos direcciones»\\
R9 & reducción al absurdo (\(\neg\)-intro) & «supongamos \(P\); contradicción; luego no \(P\)»\\
R10, R11 & contradicción / explosión & «absurdo»; «de lo falso se sigue cualquier cosa»\\
R12--R14 & principios clásicos & «por contradicción, \(P\)»; «\(P\) o no \(P\)»; «no no \(P\), luego \(P\)»\\
R15, R16 & \(\forall\)-intro / elim & «sea \(x\) arbitrario»; «en particular para \(t\)»\\
R17 & sustitución de iguales & «como \(a=b\), reemplazando»\\
\bottomrule
\end{tabular}}
\end{center}

\subsection{«Con la premisa conservada en el predicado»}

Esta expresión de la tesis (convenio 5 de la Nota 11.2) es la técnica de
\textbf{generalizar antes de inducir}. Si una hipótesis cambia entre el caso
grande y el caso pequeño, no se fija: se mete dentro de la propiedad que se
prueba por inducción. Así la hipótesis inductiva puede usarse con otra
instancia de esa hipótesis. Aparece en casi todos los lemas del capítulo:
\begin{center}
{\small
\begin{tabular}{@{}lll@{}}
\toprule
\textbf{Lema} & \textbf{Lo que se generaliza} & \textbf{Por qué}\\
\midrule
11.9 debilitamiento & el contexto mayor \(\Delta\) & las reglas con descarga piden \(\Delta\cup\{\alpha\}\)\\
11.15 Kalmár & la cobertura \(\atomos(\theta)\subseteq A\) & cada subfórmula tiene su propia cobertura\\
11.17 descarga & «\(\forall v\)» y \(\Nodup\) & la cola se usa con valuaciones modificadas\\
A.3 corrección & la valuación \(v\) y \(v\models\Gamma\) & la premisa de \(\Imp\)I tiene contexto \(\Gamma\cup\{\alpha\}\)\\
\bottomrule
\end{tabular}}
\end{center}

%=====================================================================
\section{Los resultados, uno por uno}
%=====================================================================

\subsection{Reglas estructurales y derivadas}

%---------------------------------------------------------------------
\begin{resultado}{Lema 11.9 · Debilitamiento}
\enunciado Si \(\Gamma\Provnd\varphi\) y \(\Gamma\subseteq\Delta\), entonces
\(\Delta\Provnd\varphi\).

\palabras Añadir hipótesis nunca impide una derivación.

\sirve Para «trasladar» una derivación ya hecha a un contexto con más
hipótesis. Es el paso marcado \emph{debil.} en todos los árboles del capítulo:
en los Lemas 11.10, 11.11 y 11.15.

\idea Se rehace la misma derivación, regla por regla, con \(\Delta\) en lugar
de \(\Gamma\).

\prueba Por inducción sobre la derivación de \(\Gamma\Provnd\varphi\). Lo que
se demuestra es la propiedad
\[
P(\Gamma,\varphi):\qquad \text{para todo }\Delta\supseteq\Gamma,\quad \Delta\Provnd\varphi .
\]
Un hecho auxiliar: si \(\Gamma\subseteq\Delta\), entonces
\(\Gamma\cup\{\alpha\}\subseteq\Delta\cup\{\alpha\}\).

\begin{center}
{\small
\renewcommand{\arraystretch}{1.25}
\begin{tabular}{@{}lll@{}}
\toprule
\textbf{Última regla} & \textbf{Hipótesis inductiva usada en} & \textbf{Se concluye con}\\
\midrule
hip.\ (\(\varphi\in\Gamma\)) & --- & \(\varphi\in\Gamma\subseteq\Delta\), regla hip.\\
\(\Imp\)E, \(\Not\)E, \(\bot\)E & el mismo \(\Delta\) & la misma regla\\
\(\Imp\)I, \(\Not\)I (premisa \(\Gamma,\alpha\Provnd\cdot\)) & \(\Delta\cup\{\alpha\}\) & la misma regla\\
RAA (premisa \(\Gamma,\Not\alpha\Provnd\bot\)) & \(\Delta\cup\{\Not\alpha\}\) & RAA\\
\bottomrule
\end{tabular}}
\end{center}
En los tres casos con descarga, el hecho auxiliar da
\(\Gamma\cup\{\alpha\}\subseteq\Delta\cup\{\alpha\}\), y la hipótesis
inductiva se aplica a ese contexto mayor.
\end{resultado}

\begin{trampa}
\begin{itemize}
  \item \emph{¿Por qué hay que probarlo si los contextos son conjuntos?}
    Porque los conjuntos solo hacen gratuitos el orden y las repeticiones.
    Que una derivación siga valiendo con más hipótesis es un hecho sobre las
    \emph{derivaciones}, y se demuestra recorriéndolas.
  \item \emph{¿Por qué \(\Delta\) tiene que estar cuantificado?} En el caso
    \(\Imp\)I la premisa es \(\Gamma\cup\{\alpha\}\Provnd\beta\), y hay que
    aplicarle la hipótesis inductiva con \(\Delta\cup\{\alpha\}\), no con
    \(\Delta\). Si \(\Delta\) estuviera fijo, la hipótesis inductiva no
    serviría.
\end{itemize}
\end{trampa}

%---------------------------------------------------------------------
\begin{resultado}{Lema 11.10 · Introducción de doble negación}
\enunciado Si \(\Gamma\Provnd\varphi\), entonces \(\Gamma\Provnd\Not\Not\varphi\).

\palabras Lo que se demuestra también se demuestra doblemente negado.

\sirve En el lema de Kalmár, caso \(\Not p\) con fila falsa: hay que
derivar \(\Not(\Not p)\) a partir de \(p\).

\idea Se supone \(\Not\varphi\); junto con \(\varphi\) da \(\bot\); se
descarga.

\prueba
\begin{center}
\AxiomC{}\RightLabel{\scriptsize hip.}\UnaryInfC{\(\Gamma,\Not\varphi\Provnd\Not\varphi\)}
\AxiomC{\(\Gamma\Provnd\varphi\)}\RightLabel{\scriptsize debil.}\UnaryInfC{\(\Gamma,\Not\varphi\Provnd\varphi\)}
\RightLabel{\scriptsize\(\Not\)E}\BinaryInfC{\(\Gamma,\Not\varphi\Provnd\bot\)}
\RightLabel{\scriptsize\(\Not\)I}\UnaryInfC{\(\Gamma\Provnd\Not\Not\varphi\)}
\DisplayProof
\end{center}
Se usan las reglas hip., \(\Not\)E y \(\Not\)I, más el debilitamiento con
\(\Gamma\subseteq\Gamma\cup\{\Not\varphi\}\).
\end{resultado}

\begin{trampa}
Usa \(\Not\)I y no RAA, así que es intuicionista. La dirección contraria,
\(\Not\Not\varphi\Rightarrow\varphi\), sí necesita RAA.
\end{trampa}

%---------------------------------------------------------------------
\begin{resultado}{Lema 11.11 · Análisis por casos derivado}
\enunciado Si \(\Gamma,\varphi\Provnd\chi\) y \(\Gamma,\Not\varphi\Provnd\chi\),
entonces \(\Gamma\Provnd\chi\).

\palabras Si \(\chi\) se sigue tanto suponiendo \(\varphi\) como suponiendo
\(\Not\varphi\), entonces se sigue sin suponer nada sobre \(\varphi\). Es el
tercero excluido, pero escrito sin el conectivo \(\vee\).

\sirve Es el motor de la descarga (Lema 11.17): elimina un literal del
contexto.

\idea Por RAA: se supone \(\Not\chi\) y se demuestra que tanto \(\varphi\)
como \(\Not\varphi\) llevan a contradicción.

\prueba Sea \(\Delta:=\Gamma\cup\{\Not\chi\}\) (se abre \(\Not\chi\) para RAA).
\begin{enumerate}
  \item \(\Delta\Provnd\Not\varphi\): se supone \(\varphi\). En
    \(\Delta,\varphi\) se tiene \(\Not\chi\) (hip.) y \(\chi\) (primera premisa,
    debilitada); \(\Not\)E da \(\bot\) y \(\Not\)I descarga \(\varphi\).
  \item \(\Delta\Provnd\Not\Not\varphi\): igual, suponiendo \(\Not\varphi\) y
    usando la segunda premisa.
  \item \(\Not\)E con 1 y 2 da \(\Delta\Provnd\bot\), es decir,
    \(\Gamma,\Not\chi\Provnd\bot\).
  \item RAA descarga \(\Not\chi\) y da \(\Gamma\Provnd\chi\).
\end{enumerate}
\begin{center}
\AxiomC{}\RightLabel{\scriptsize hip.}\UnaryInfC{\(\Delta,\varphi\Provnd\Not\chi\)}
\AxiomC{\(\Gamma,\varphi\Provnd\chi\)}\RightLabel{\scriptsize debil.}\UnaryInfC{\(\Delta,\varphi\Provnd\chi\)}
\RightLabel{\scriptsize\(\Not\)E}\BinaryInfC{\(\Delta,\varphi\Provnd\bot\)}
\RightLabel{\scriptsize\(\Not\)I}\UnaryInfC{\(\Delta\Provnd\Not\varphi\)}
\DisplayProof
\end{center}
(El árbol de \(\Delta\Provnd\Not\Not\varphi\) es idéntico, cambiando
\(\varphi\) por \(\Not\varphi\).)
\end{resultado}

\begin{trampa}
Es el \textbf{único lugar} en que la prueba de completitud usa la regla
clásica RAA \emph{dentro del cálculo}. Las derivaciones del lema de Kalmár, el
debilitamiento y la doble negación usan solo reglas intuicionistas. Esta
observación se lee directamente en las pruebas de la tesis. Explica por qué la
completitud falla sin RAA: sin este lema no se pueden juntar las filas.
\end{trampa}

\subsection{El lema de Kalmár y sus auxiliares}

%---------------------------------------------------------------------
\begin{resultado}{Lema 11.14 · Pertenencia del literal}
\enunciado Si \(p\in A\), entonces \(\lit_v(p)\in\Gamma_v(A)\).

\palabras Si el átomo está en la lista, su literal está en el contexto.

\sirve En el caso atómico de Kalmár, para cerrar con la regla de hipótesis.

\prueba Por inducción sobre la lista \(A\) (con la premisa \(p\in A\) dentro
de la propiedad).
\begin{itemize}
  \item \(A=[\,]\): \(p\in[\,]\) es imposible; el caso es vacío.
  \item \(A=a::r\): o bien \(p=a\), y entonces \(\lit_v(p)=\lit_v(a)\) es el
    elemento que se añadió; o bien \(p\in r\), y por hipótesis inductiva
    \(\lit_v(p)\in\Gamma_v(r)\subseteq\Gamma_v(a::r)\).
\end{itemize}
\end{resultado}

%---------------------------------------------------------------------
\begin{resultado}{Lema 11.15 · Lema de Kalmár}
\enunciado Sea \(A\) una lista que cubre a \(\varphi\). Para toda valuación
\(v\),
\[
\sem\varphi_v\ \Rightarrow\ \Gamma_v(A)\Provnd\varphi
\qquad\text{y}\qquad
\text{no }\sem\varphi_v\ \Rightarrow\ \Gamma_v(A)\Provnd\Not\varphi .
\]

\palabras En cada fila de la tabla, los literales de esa fila demuestran la
fórmula si es verdadera, y demuestran su negación si es falsa.

\sirve Es la primera etapa de la completitud: da una derivación por fila.

\idea Inducción sobre \(\varphi\). En cada caso se mira qué dice la tabla de
verdad del conectivo y se construye la derivación correspondiente.

\prueba Sea \(\Gamma:=\Gamma_v(A)\). Se prueba por inducción estructural la
propiedad
\[
K(\theta):\quad \atomos(\theta)\subseteq A\ \Rightarrow\
\big(\sem\theta_v\Rightarrow\Gamma\Provnd\theta\big)\ \wedge\
\big(\text{no }\sem\theta_v\Rightarrow\Gamma\Provnd\Not\theta\big).
\]
La cobertura pasa a las subfórmulas, porque sus átomos son parte de los de
\(\theta\).
\end{resultado}

\begin{center}
\renewcommand{\arraystretch}{1.45}
{\small
\begin{tabularx}{\linewidth}{@{}p{0.13\linewidth}p{0.17\linewidth}X@{}}
\toprule
\textbf{Caso} & \textbf{Fila} & \textbf{Derivación que se construye}\\
\midrule
\(\theta=s\) & \(v(s)=\V\) & \(\lit_v(s)=s\in\Gamma\) (Lema 11.14); regla hip.\\
& \(v(s)=\F\) & \(\lit_v(s)=\Not s\in\Gamma\); regla hip.\\
\midrule
\(\theta=\Not p\) & \(\sem{\Not p}_v\) verdadera & entonces \(p\) es falsa; la hipótesis inductiva negativa da \(\Gamma\Provnd\Not p\).\\
& \(\sem{\Not p}_v\) falsa & entonces \(p\) es verdadera (doble negación meta); la hipótesis inductiva positiva da \(\Gamma\Provnd p\), y el Lema 11.10 da \(\Gamma\Provnd\Not\Not p\).\\
\midrule
\(\theta=p\Imp q\) & \(p\) verdadera, \(q\) verdadera & \(\Gamma\Provnd q\); se debilita a \(\Gamma,p\Provnd q\); \(\Imp\)I.\\
& \(p\) falsa & \(\Gamma\Provnd\Not p\); en \(\Gamma,p\): \(\Not p\) y \(p\) dan \(\bot\), \(\bot\)E da \(q\); \(\Imp\)I.\\
& \(p\) verdadera, \(q\) falsa & \(\Gamma\Provnd p\) y \(\Gamma\Provnd\Not q\); en \(\Gamma,(p\Imp q)\): \(\Imp\)E da \(q\), con \(\Not q\) da \(\bot\); \(\Not\)I da \(\Gamma\Provnd\Not(p\Imp q)\).\\
\bottomrule
\end{tabularx}}
\end{center}

Las tres derivaciones del caso implicación, tal como aparecen en la tesis:
\begin{center}
\AxiomC{\(\Gamma\Provnd q\)}\RightLabel{\scriptsize debil.}\UnaryInfC{\(\Gamma,p\Provnd q\)}
\RightLabel{\scriptsize\(\Imp\)I}\UnaryInfC{\(\Gamma\Provnd p\Imp q\)}\DisplayProof
\qquad
\AxiomC{\(\Gamma\Provnd\Not p\)}\RightLabel{\scriptsize debil.}\UnaryInfC{\(\Gamma,p\Provnd\Not p\)}
\AxiomC{}\RightLabel{\scriptsize hip.}\UnaryInfC{\(\Gamma,p\Provnd p\)}
\RightLabel{\scriptsize\(\Not\)E}\BinaryInfC{\(\Gamma,p\Provnd\bot\)}
\RightLabel{\scriptsize\(\bot\)E}\UnaryInfC{\(\Gamma,p\Provnd q\)}
\RightLabel{\scriptsize\(\Imp\)I}\UnaryInfC{\(\Gamma\Provnd p\Imp q\)}\DisplayProof
\end{center}
\begin{center}
\AxiomC{\(\Gamma\Provnd\Not q\)}\RightLabel{\scriptsize debil.}\UnaryInfC{\(\Gamma,(p\Imp q)\Provnd\Not q\)}
\AxiomC{}\RightLabel{\scriptsize hip.}\UnaryInfC{\(\Gamma,(p\Imp q)\Provnd p\Imp q\)}
\AxiomC{\(\Gamma\Provnd p\)}\RightLabel{\scriptsize debil.}\UnaryInfC{\(\Gamma,(p\Imp q)\Provnd p\)}
\RightLabel{\scriptsize\(\Imp\)E}\BinaryInfC{\(\Gamma,(p\Imp q)\Provnd q\)}
\RightLabel{\scriptsize\(\Not\)E}\BinaryInfC{\(\Gamma,(p\Imp q)\Provnd\bot\)}
\RightLabel{\scriptsize\(\Not\)I}\UnaryInfC{\(\Gamma\Provnd\Not(p\Imp q)\)}\DisplayProof
\end{center}

\begin{trampa}
\begin{itemize}
  \item \emph{¿Por qué el enunciado es bilateral?} Porque en los casos
    \(\Not\) e \(\Imp\) la verdad de la fórmula depende de que una subfórmula
    sea \emph{falsa}. Para derivar \(\Not p\) cuando \(p\) es falsa, o
    \(p\Imp q\) cuando \(p\) es falsa, hace falta que la hipótesis inductiva
    diga algo también de las fórmulas falsas. Un enunciado unilateral no
    pasaría la inducción.
  \item \emph{¿Por qué «\(A\) cubre a \(\varphi\)» y no «\(A=\atomos(\varphi)\)»?}
    Porque las subfórmulas tienen menos átomos. Con cobertura, la misma lista
    \(A\) sirve en toda la inducción. Los literales sobrantes quedan en el
    contexto sin usarse: la regla hip.\ solo exige pertenencia.
  \item \emph{¿Es lógica clásica?} Las derivaciones que se construyen son
    intuicionistas (no usan RAA). Lo clásico está en el razonamiento
    \emph{sobre} las filas: «\(p\) es verdadera o falsa», «no no \(p\), luego
    \(p\)», y obtener «\(p\) verdadera y \(q\) falsa» de «no
    (\(p\Rightarrow q\))».
\end{itemize}
\end{trampa}

%---------------------------------------------------------------------
\begin{resultado}{Lema 11.16 · Congruencia de literales y contextos}
\enunciado Si \(v(x)\iff w(x)\), entonces \(\lit_v(x)=\lit_w(x)\). Por tanto,
si \(v\) y \(w\) coinciden en todos los átomos de \(A\), entonces
\(\Gamma_v(A)=\Gamma_w(A)\).

\palabras El contexto de literales solo depende de los valores de los átomos
de la lista.

\sirve En la descarga, al cambiar el valor de un solo átomo \(p\), garantiza
que el resto del contexto no cambia.

\prueba Para los literales, por casos sobre \(v(x)\): si es verdadero, también
lo es \(w(x)\), y ambos literales son \(x\); si es falso, \(w(x)\) también lo
es, y ambos son \(\Not x\). Para los contextos, por inducción sobre \(A\),
usando lo anterior en la cabeza y la hipótesis inductiva en la cola.
\end{resultado}

%---------------------------------------------------------------------
\begin{resultado}{Lema 11.17 · Descarga de los átomos}
\enunciado Sea \(A\) una lista sin repeticiones. Si
\(\Gamma_v(A)\Provnd\varphi\) para \emph{toda} valuación \(v\), entonces
\(\varnothing\Provnd\varphi\).

\palabras Si la fórmula se deriva en todas las filas, se deriva sin
hipótesis.

\sirve Es la segunda etapa de la completitud.

\idea Se eliminan los átomos de la lista uno a uno, empezando por la cabeza.
Para eliminar \(p\) se juntan, con el Lema 11.11, las dos filas que solo
difieren en \(p\).

\prueba Por inducción sobre \(A\), demostrando
\[
D(s):\quad \Nodup(s)\ \Rightarrow\ \big(\forall v,\ \Gamma_v(s)\Provnd\varphi\big)\ \Rightarrow\ \varnothing\Provnd\varphi .
\]
\begin{itemize}
  \item \textbf{\(A=[\,]\).} \(\Gamma_v([\,])=\varnothing\) para cualquier
    \(v\). Basta tomar una valuación cualquiera (la tesis usa la que hace
    todo verdadero).
  \item \textbf{\(A=p::r\).} Por \(\Nodup(p::r)\) se tiene \(p\notin r\) y
    \(\Nodup(r)\). Para aplicar la hipótesis inductiva hay que probar
    \(\forall v,\ \Gamma_v(r)\Provnd\varphi\). Sea \(v\) cualquiera y sean
    \[
    u:=v\text{ con }p\mapsto\V,\qquad w:=v\text{ con }p\mapsto\F .
    \]
    \begin{enumerate}
      \item Como \(p\notin r\), \(u\), \(v\) y \(w\) coinciden en los átomos
        de \(r\); por el Lema 11.16,
        \(\Gamma_u(r)=\Gamma_v(r)=\Gamma_w(r)=:\Delta\).
      \item Entonces \(\Gamma_u(p::r)=\Delta\cup\{p\}\) y
        \(\Gamma_w(p::r)=\Delta\cup\{\Not p\}\).
      \item La hipótesis «para toda valuación», aplicada a \(u\) y a \(w\),
        da \(\Delta,p\Provnd\varphi\) y \(\Delta,\Not p\Provnd\varphi\).
      \item El Lema 11.11 da \(\Delta\Provnd\varphi\), es decir,
        \(\Gamma_v(r)\Provnd\varphi\).
    \end{enumerate}
    Como \(v\) era arbitraria, se tiene \(\forall v,\ \Gamma_v(r)\Provnd\varphi\),
    y la hipótesis inductiva da \(\varnothing\Provnd\varphi\).
\end{itemize}
\end{resultado}

\begin{center}
\begin{tikzpicture}[font=\footnotesize, level distance=15mm,
  level 1/.style={sibling distance=60mm}, level 2/.style={sibling distance=30mm},
  every node/.style={draw, rounded corners, fill=blue!6, inner sep=3pt},
  lab/.style={draw=none, fill=none, font=\scriptsize}]
\node[fill=yellow!20] {$\varnothing\Provnd\varphi$}
  child { node {$\{p\}\Provnd\varphi$}
    child { node[fill=green!12] {$\{q,p\}\Provnd\varphi$} edge from parent node[lab,left] {$q$} }
    child { node[fill=green!12] {$\{\Not q,p\}\Provnd\varphi$} edge from parent node[lab,right] {$\Not q$} }
    edge from parent node[lab,left] {$p$} }
  child { node {$\{\Not p\}\Provnd\varphi$}
    child { node[fill=green!12] {$\{q,\Not p\}\Provnd\varphi$} edge from parent node[lab,left] {$q$} }
    child { node[fill=green!12] {$\{\Not q,\Not p\}\Provnd\varphi$} edge from parent node[lab,right] {$\Not q$} }
    edge from parent node[lab,right] {$\Not p$} };
\end{tikzpicture}
\end{center}
{\small Descarga con \(A=[q,p]\). Las hojas (verdes) son las cuatro
instancias del lema de Kalmár. Cada nodo interno es una aplicación del
Lema 11.11. Primero se descarga la cabeza \(q\) (nivel inferior) y después
\(p\). Con \(k\) átomos hay \(2^k\) hojas y \(2^k-1\) aplicaciones del
análisis por casos.}

\begin{trampa}
\begin{itemize}
  \item \emph{¿Para qué la hipótesis «para toda \(v\)»?} Porque para eliminar
    \(p\) se necesitan las dos filas \(u\) y \(w\), que no son la fila \(v\) de
    partida. Por eso «\(\forall v\)» va dentro de la propiedad que se prueba
    por inducción.
  \item \emph{¿Para qué \(\Nodup\)?} Garantiza que cambiar el valor de \(p\) no
    altera los literales de la cola (paso 1), de modo que los contextos
    coinciden \emph{exactamente} y las premisas encajan con el Lema 11.11.
    Observación propia, no de la tesis: con repeticiones se podría argumentar
    con inclusiones y debilitamiento, pero la prueba sería menos directa.
    Esta es la razón práctica de \(\dedup\) en el Teorema 11.18.
  \item \emph{¿Por qué no se usan conjuntos maximales (Lindenbaum--Henkin)?}
    Porque una fórmula tiene un número finito de átomos: basta recorrer la
    tabla de verdad, que es finita. El método de Kalmár es constructivo en el
    sentido de que produce la derivación fila por fila.
\end{itemize}
\end{trampa}

%---------------------------------------------------------------------
\begin{resultado}{Teorema 11.18 · Completitud}
\enunciado \(\Taut(\varphi)\Rightarrow\varnothing\Provnd\varphi\).

\prueba Supóngase \(\Taut(\varphi)\) y sea \(A:=\dedup(\atomos(\varphi))\).
\begin{enumerate}
  \item \(A\) no tiene repeticiones y contiene los mismos átomos que
    \(\atomos(\varphi)\); en particular, cubre a \(\varphi\).
  \item Sea \(v\) cualquiera. Como \(\varphi\) es tautología,
    \(\sem\varphi_v\) es verdadera, y la parte positiva del lema de Kalmár da
    \(\Gamma_v(A)\Provnd\varphi\).
  \item Así, \(\forall v,\ \Gamma_v(A)\Provnd\varphi\). La descarga
    (Lema 11.17) da \(\varnothing\Provnd\varphi\).
\end{enumerate}
\end{resultado}

\begin{nota}[Por qué solo se usa la parte positiva de Kalmár]
En el teorema final todas las filas son verdaderas, así que solo se usa la
mitad positiva. La mitad negativa es imprescindible \emph{dentro} de la
inducción del lema (casos \(\Not\) e \(\Imp\)), no fuera de ella.
\end{nota}

\subsection{Apéndice A: corrección, consistencia y adecuación}

%---------------------------------------------------------------------
\begin{resultado}{Lema A.1 · \(\bot\) es falsa en toda valuación}
\enunciado \(\sem{\bot}_v\iff\text{Falso}\).

\prueba \(\sem\bot_v=\text{no}\,(v(p_0)\Rightarrow v(p_0))\), y
\(v(p_0)\Rightarrow v(p_0)\) es siempre verdadera. Luego \(\sem\bot_v\) es
falsa. La otra dirección es la explosión: de Falso se sigue cualquier cosa.
\end{resultado}

%---------------------------------------------------------------------
\begin{resultado}{Lema A.2 · Satisfacción de un contexto extendido}
\enunciado Si \(\sem\chi_v\) y \(v\models\Gamma\), entonces
\(v\models\Gamma\cup\{\chi\}\).

\prueba Toda \(\psi\in\Gamma\cup\{\chi\}\) es \(\chi\) (verdadera por
hipótesis) o está en \(\Gamma\) (verdadera porque \(v\models\Gamma\)).
\end{resultado}

%---------------------------------------------------------------------
\begin{resultado}{Teorema A.3 · Corrección}
\enunciado Si \(\Gamma\Provnd\varphi\), entonces para toda \(v\) con
\(v\models\Gamma\) se cumple \(\sem\varphi_v\).

\palabras Las reglas solo producen conclusiones verdaderas a partir de
hipótesis verdaderas.

\idea Cada regla preserva la verdad.

\prueba Por inducción sobre la derivación, con la valuación cuantificada
dentro de la propiedad. Sea \(v\models\Gamma\).
\begin{center}
{\small
\renewcommand{\arraystretch}{1.3}
\begin{tabularx}{\linewidth}{@{}p{0.13\linewidth}X@{}}
\toprule
\textbf{Regla} & \textbf{Por qué la conclusión es verdadera}\\
\midrule
hip. & \(\varphi\in\Gamma\) y \(v\models\Gamma\).\\
\(\Imp\)I & Si \(\sem\alpha_v\), entonces \(v\models\Gamma\cup\{\alpha\}\) (A.2) y la hipótesis inductiva da \(\sem\beta_v\). Luego \(\sem{\alpha\Imp\beta}_v\).\\
\(\Imp\)E & Las hipótesis inductivas dan \(\sem\alpha_v\Rightarrow\sem\beta_v\) y \(\sem\alpha_v\).\\
\(\Not\)I & Si \(\sem\alpha_v\), la hipótesis inductiva daría \(\sem\bot_v\), imposible (A.1). Luego no \(\sem\alpha_v\).\\
\(\Not\)E & Darían \(\sem\alpha_v\) y no \(\sem\alpha_v\): contradicción, de la que se sigue \(\sem\bot_v\).\\
\(\bot\)E & La hipótesis inductiva da \(\sem\bot_v\), imposible; de lo falso se sigue \(\sem\alpha_v\).\\
RAA & Si \(\sem\alpha_v\) fuera falsa, \(v\models\Gamma\cup\{\Not\alpha\}\) y se tendría \(\sem\bot_v\), imposible. \textbf{Luego} \(\sem\alpha_v\) es verdadera.\\
\bottomrule
\end{tabularx}}
\end{center}
\end{resultado}

\begin{trampa}
El caso RAA es el único que razona por contradicción en el nivel meta: pasa de
«no puede ser falsa» a «es verdadera». Eso es la \emph{bivalencia} de la
semántica de tablas. Es la razón exacta por la que RAA es sólida para las
tablas de verdad y no para semánticas no bivalentes (Kripke).
\end{trampa}

%---------------------------------------------------------------------
\begin{resultado}{Corolarios A.4 y A.5 · Teoremas son tautologías; consistencia}
\enunciado (A.4) \(\varnothing\Provnd\varphi\Rightarrow\Taut(\varphi)\).
(A.5) \(\varnothing\not\Provnd\bot\).

\prueba (A.4) Toda \(v\) satisface \(\varnothing\) (no hay nada que
comprobar), así que la corrección da \(\sem\varphi_v\) para toda \(v\).
(A.5) Si \(\varnothing\Provnd\bot\), por A.4 \(\bot\) sería tautología;
pero es falsa en la valuación que hace todo verdadero (A.1). Contradicción.
\end{resultado}

%---------------------------------------------------------------------
\begin{resultado}{Teorema A.6 · Adecuación}
\enunciado \(\Taut(\varphi)\iff\varnothing\Provnd\varphi\).

\prueba La dirección \((\Rightarrow)\) es la completitud (11.18); la
\((\Leftarrow)\), el Corolario A.4.
\end{resultado}

%=====================================================================
\section{Ejemplo resuelto de principio a fin}
%=====================================================================

Se sigue la prueba completa con
\[
\varphi\ :=\ p\Imp(q\Imp p).
\]

\textbf{Paso 0: es tautología.} Si \(p\) es verdadera, \(q\Imp p\) también lo
es. Si \(p\) es falsa, la implicación externa es verdadera. Por tanto, las cuatro
filas dan \(\V\).

\textbf{Paso 1: la lista de átomos.} \(\atomos(\varphi)=[p]\mathbin{+\!\!+}([q]\mathbin{+\!\!+}[p])=[p,q,p]\).
Tiene una repetición; \(\dedup([p,q,p])=[q,p]=:A\). Así, \(A\) no tiene
repeticiones y cubre a \(\varphi\).

\textbf{Paso 2: Kalmár, fila por fila.} \(\Gamma_v(A)=\{\lit_v(q),\lit_v(p)\}\).
La prueba del lema mira primero \(p\) (antecedente de la implicación externa).

\begin{ejemplo}{Filas con \(v(p)=\V\): se deriva \(q\Imp p\) y se aplica \(\Imp\)I}
\emph{Fila \(p=\V,q=\V\)}, \(\Gamma=\{q,p\}\). Kalmár sobre \(q\Imp p\) con
\(q\) verdadera: \(\Gamma\Provnd p\) por hipótesis.
\begin{center}
\AxiomC{}\RightLabel{\scriptsize hip.}\UnaryInfC{\(\Gamma\Provnd p\)}
\RightLabel{\scriptsize debil.}\UnaryInfC{\(\Gamma,q\Provnd p\)}
\RightLabel{\scriptsize\(\Imp\)I}\UnaryInfC{\(\Gamma\Provnd q\Imp p\)}
\RightLabel{\scriptsize debil.}\UnaryInfC{\(\Gamma,p\Provnd q\Imp p\)}
\RightLabel{\scriptsize\(\Imp\)I}\UnaryInfC{\(\Gamma\Provnd p\Imp(q\Imp p)\)}
\DisplayProof
\end{center}
(Como los contextos son conjuntos, \(\Gamma,q=\Gamma\) y \(\Gamma,p=\Gamma\):
el debilitamiento aquí no añade nada, pero la prueba es uniforme.)

\emph{Fila \(p=\V,q=\F\)}, \(\Gamma=\{\Not q,p\}\). Kalmár sobre \(q\Imp p\)
con \(q\) falsa usa la rama «antecedente falso»:
\begin{center}
\AxiomC{}\RightLabel{\scriptsize hip.}\UnaryInfC{\(\Gamma\Provnd\Not q\)}
\RightLabel{\scriptsize debil.}\UnaryInfC{\(\Gamma,q\Provnd\Not q\)}
\AxiomC{}\RightLabel{\scriptsize hip.}\UnaryInfC{\(\Gamma,q\Provnd q\)}
\RightLabel{\scriptsize\(\Not\)E}\BinaryInfC{\(\Gamma,q\Provnd\bot\)}
\RightLabel{\scriptsize\(\bot\)E}\UnaryInfC{\(\Gamma,q\Provnd p\)}
\RightLabel{\scriptsize\(\Imp\)I}\UnaryInfC{\(\Gamma\Provnd q\Imp p\)}
\DisplayProof
\end{center}
y después debil.\ y \(\Imp\)I, como arriba.
\end{ejemplo}

\begin{ejemplo}{Filas con \(v(p)=\F\): antecedente falso}
\(\Gamma=\{\lit_v(q),\Not p\}\), con cualquier valor de \(q\).
\begin{center}
\AxiomC{}\RightLabel{\scriptsize hip.}\UnaryInfC{\(\Gamma\Provnd\Not p\)}
\RightLabel{\scriptsize debil.}\UnaryInfC{\(\Gamma,p\Provnd\Not p\)}
\AxiomC{}\RightLabel{\scriptsize hip.}\UnaryInfC{\(\Gamma,p\Provnd p\)}
\RightLabel{\scriptsize\(\Not\)E}\BinaryInfC{\(\Gamma,p\Provnd\bot\)}
\RightLabel{\scriptsize\(\bot\)E}\UnaryInfC{\(\Gamma,p\Provnd q\Imp p\)}
\RightLabel{\scriptsize\(\Imp\)I}\UnaryInfC{\(\Gamma\Provnd p\Imp(q\Imp p)\)}
\DisplayProof
\end{center}
\end{ejemplo}

\textbf{Paso 3: descarga.} Con \(A=[q,p]\) se descarga primero \(q\):
\[
\frac{\{q,p\}\Provnd\varphi\quad\{\Not q,p\}\Provnd\varphi}{\{p\}\Provnd\varphi}\ \text{(11.11)}
\qquad
\frac{\{q,\Not p\}\Provnd\varphi\quad\{\Not q,\Not p\}\Provnd\varphi}{\{\Not p\}\Provnd\varphi}\ \text{(11.11)}
\]
y después \(p\):
\[
\frac{\{p\}\Provnd\varphi\qquad\{\Not p\}\Provnd\varphi}{\varnothing\Provnd p\Imp(q\Imp p)}\ \text{(11.11)}.
\]
Es exactamente el árbol dibujado tras el Lema 11.17: cuatro hojas y tres
aplicaciones del análisis por casos. Cada aplicación usa una vez RAA.

\begin{nota}[Qué muestra el ejemplo]
La derivación final no es la más corta (la directa usa dos \(\Imp\)I y una
hipótesis). El método de Kalmár es \emph{uniforme}: no busca la mejor prueba,
sino que garantiza que \emph{alguna} existe, con un procedimiento que funciona
para cualquier tautología.
\end{nota}

%=====================================================================
\section{Dónde interviene la lógica clásica}
%=====================================================================

\begin{center}
\renewcommand{\arraystretch}{1.35}
{\small
\begin{tabularx}{\linewidth}{@{}p{0.27\linewidth}XX@{}}
\toprule
\textbf{Resultado} & \textbf{Dentro del cálculo \(\Provnd\)} & \textbf{En el razonamiento sobre el cálculo}\\
\midrule
11.9, 11.10, 11.14, 11.16 & sin RAA & ---\\
11.11 análisis por casos & \textbf{RAA} (única regla clásica usada) & ---\\
11.15 Kalmár & sin RAA & \(p\) verdadera o falsa; doble negación; de «no \((P\Rightarrow Q)\)» a \(P\)\\
11.17 descarga & RAA, a través de 11.11 & ---\\
Definición de literal & --- & decidir si \(v(p)\) es verdadera\\
A.3 corrección & --- & caso RAA: «no es falsa, luego es verdadera»\\
\bottomrule
\end{tabularx}}
\end{center}

Para resumirlo en la defensa: \emph{en el cálculo}, lo clásico entra una sola
vez, al juntar filas (Lema 11.11). \emph{En el metanivel}, entra porque la
semántica es de tablas de verdad, es decir, bivalente.

\subsection*{Qué garantiza la auditoría de axiomas (§11.6)}

La herramienta lista los axiomas lógicos de los que depende el teorema. El
resultado son tres axiomas estándar de la matemática clásica en el sistema
(extensionalidad proposicional, elección y solidez de los cocientes), sin
ninguna demostración pendiente ni una «regla de completitud» postulada.
\begin{itemize}
  \item \textbf{Lo que sí garantiza:} no hay huecos admitidos y no se postuló
    el resultado.
  \item \textbf{Lo que no garantiza por sí sola:} que no se haya usado un
    teorema de completitud de una biblioteca. Esa independencia se justifica
    por la estructura de los archivos y por las pruebas (§11.6).
\end{itemize}
Lo mismo vale para la corrección y la adecuación (§A.3).

%=====================================================================
\section{Preguntas probables del tribunal}
%=====================================================================

\begin{description}[style=nextline, leftmargin=0pt, itemsep=6pt]
\item[¿Qué demuestra exactamente el capítulo?]
Que toda tautología del fragmento \(\{\Not,\Imp\}\) es derivable sin
hipótesis. Con el Apéndice A, que ser tautología y ser derivable sin
hipótesis son equivalentes.

\item[¿Por qué solo para el contexto vacío?]
Porque es el caso que se formaliza. Para un \(\Gamma=\{\gamma_1,\dots,\gamma_n\}\)
finito con \(\Gamma\models\psi\), la fórmula
\(\gamma_1\Imp(\cdots\Imp(\gamma_n\Imp\psi))\) es tautología; por completitud
es derivable sin hipótesis; con debilitamiento y \(n\) aplicaciones de
\(\Imp\)E se obtiene \(\Gamma\Provnd\psi\). La tesis describe este paso en el
esbozo del Capítulo 9; no está formalizado.

\item[¿Por qué la base \(\{\Not,\Imp\}\)?]
Es funcionalmente completa: toda tabla de verdad se expresa con ella. Además
reduce la inducción a tres casos. La traducción \(T\) preserva la verdad
(§\ref{sec:base}).

\item[¿Por qué \(\bot\) es una abreviatura y \(\Not\) es primitiva?]
Porque el lenguaje no tiene \(\bot\). Se define \(\bot:=\Not(p_0\Imp p_0)\),
falsa en toda valuación. Mendelson y Stansifer usan la misma base
\(\{\Not,\Imp\}\); la falsedad se codifica dentro del fragmento.

\item[¿Por qué hay que demostrar el debilitamiento?]
Los conjuntos dan gratis el orden y las repeticiones, pero no que una
derivación siga valiendo con más hipótesis. Eso es un hecho sobre las
derivaciones.

\item[¿Por qué el lema de Kalmár es bilateral?]
Los casos \(\Not\) e \(\Imp\) dependen de subfórmulas falsas. Sin la parte
negativa la inducción no pasa.

\item[¿Qué es una cobertura y por qué no basta \(A=\atomos(\varphi)\)?]
Las subfórmulas tienen menos átomos. Con «\(A\) contiene los átomos» se usa
la misma lista en toda la inducción; los literales sobrantes no estorban.

\item[¿Para qué sirven \(\Nodup\) y \(\dedup\)?]
Para que cambiar el valor de un átomo no altere los literales de la cola, y
los contextos coincidan exactamente al aplicar el análisis por casos.

\item[¿Dónde se usa la lógica clásica?]
En el cálculo, solo en el Lema 11.11 (RAA al juntar filas). En el
razonamiento meta, por la bivalencia de las tablas de verdad.

\item[¿Por qué no Lindenbaum--Henkin?]
No hace falta: una fórmula tiene finitos átomos y la tabla de verdad es
finita. Kalmár recorre esa tabla.

\item[¿La prueba construye la derivación?]
Matemáticamente sí: el argumento es un procedimiento. En la formalización la
derivación no se exporta como dato, porque la elección del literal decide una
proposición arbitraria (§11.1).

\item[¿Cuántos pasos de análisis por casos hacen falta?]
Con \(k\) átomos distintos, \(2^k\) filas y \(2^k-1\) aplicaciones del
Lema 11.11.

\item[¿Por qué en la corrección la valuación va cuantificada dentro de la inducción?]
Por la misma razón que \(\Delta\) en el debilitamiento: la premisa de
\(\Imp\)I tiene contexto \(\Gamma\cup\{\alpha\}\), y la hipótesis inductiva
debe poder usarse con la satisfacción de ese contexto.

\item[¿Qué relación hay con el Capítulo 9?]
El Teorema 9.14 dio el esbozo; aquí está la prueba completa. El Ejemplo 9.17
(\(p\Imp p\)) es el caso de un átomo del mismo procedimiento.
\end{description}

%=====================================================================
\appendix
\section{Hoja de referencia}
%=====================================================================

\begin{clave}
\textbf{La cadena completa en una línea:}
\[
\Taut(\varphi)
\ \xrightarrow{\ \text{Kalmár (11.15) en cada fila}\ }\
\forall v,\ \Gamma_v(A)\Provnd\varphi
\ \xrightarrow{\ \text{descarga (11.17), usando 11.11}\ }\
\varnothing\Provnd\varphi .
\]
\textbf{Y de vuelta:}\quad
\(\varnothing\Provnd\varphi\xrightarrow{\ \text{corrección (A.3)}\ }\Taut(\varphi)\).
\end{clave}

\textbf{Las siete reglas:} hip., \(\Imp\)I, \(\Imp\)E, \(\Not\)I, \(\Not\)E,
\(\bot\)E (núcleo intuicionista) y RAA (clásica).

\textbf{Tres técnicas del capítulo:}
\begin{enumerate}
  \item \emph{Inducción sobre derivaciones, generalizando lo que cambia}
    (11.9, A.3).
  \item \emph{Inducción estructural con un enunciado reforzado} (11.15).
  \item \emph{Inducción sobre una lista, modificando la valuación en la
    cabeza} (11.17).
\end{enumerate}

\textbf{Los tres datos que más se preguntan:}
\begin{itemize}
  \item Kalmár es bilateral porque \(\Not\) e \(\Imp\) dependen de
    subfórmulas falsas.
  \item El único uso de RAA en el cálculo está en el análisis por casos
    (11.11).
  \item La completitud es para \(\Gamma=\varnothing\); la corrección, para
    cualquier \(\Gamma\).
\end{itemize}

\end{document}
